\section*{2026 年 4 月 27 日作业}
\phantomsection
\addcontentsline{toc}{section}{2026 年 4 月 27 日作业}
\noindent\textbf{截止时间:2026 年 5 月 11 日 13:30}
\setcounter{problem}{0}
\setcounter{sollemma}{0}

\begin{problem}
已知 $m$ 和 $n$ 是正整数, $A\in\C^{m\times n}$, $B\in\C^{n\times m}$, 证明: $AB$ 和 $BA$ 具有完全相同的非零特征值(计代数重数). 当 $m=n$ 时 $AB$ 和 $BA$ 是否一定相似?
\begin{solution}

\end{solution}
\end{problem}

\begin{problem}
已知 $n$ 是正整数, $u,v\in\C^n$, 求 $uv^*$ 的特征多项式和极小多项式.
\begin{solution}

\end{solution}
\end{problem}

\begin{problem}
已知 $n$ 是正整数. 若 $A,B\in\C^{n\times n}$ 可交换, 即 $AB=BA$, 证明: 存在酉矩阵 $Q$ 使得 $Q^*AQ$ 和 $Q^*BQ$ 都是上三角矩阵.
\begin{solution}

\end{solution}
\end{problem}

\begin{problem}
若 $(A_k)_{k\geq 1}$ 是收敛于 $A$ 的复方阵序列, $A_k=Q_kT_kQ_k^*$ 是 $A_k$ 的 Schur 分解. 证明: 在 $(Q_k)_{k\geq 1}$ 中存在收敛子列 $(Q_{k_n})_{n\geq 1}$, 使得 $Q=\lim_{n\to\infty}Q_{k_n}$ 是酉矩阵且 $Q^*AQ$ 是上三角矩阵.
\begin{solution}

\end{solution}
\end{problem}

\begin{problem}
已知 $n$ 是正整数, $A\in\R^{n\times n}$. 证明: 存在正交矩阵 $Q$ 使得 $Q^TAQ$ 是分块上三角矩阵, 且每个对角块的阶数不超过 $2$.
\begin{solution}

\end{solution}
\end{problem}

\begin{problem}[选做]
已知 $n$ 是正整数, 求 $n$ 阶三对角矩阵
\[
    A=
    \begin{bmatrix}
        0 & 1 &  &  &  \\
        n-1 & 0 & 2 &  &  \\
         & n-2 & 0 & \ddots &  \\
         &  & \ddots & \ddots & n-1 \\
         &  &  & 1 & 0
    \end{bmatrix}
\]
的所有特征值.
\begin{solution}

\end{solution}
\end{problem}

\begin{problem}[选做]
已知 $m,n$ 是正整数. 证明循环 Schur 分解: 给定 $n$ 阶复矩阵 $A_1,A_2,\ldots,A_m$, 一定存在 $n$ 阶酉矩阵 $Q_1,Q_2,\ldots,Q_m$ 使得 $Q_1^*A_1Q_2$, $Q_2^*A_2Q_3$, $\ldots$, $Q_{m-1}^*A_{m-1}Q_m$, $Q_m^*A_mQ_1$ 都是上三角矩阵.
\begin{solution}

\end{solution}
\end{problem}

\begin{problem}[选做]
若递归数列 $(\alpha_n)_{n\geq 0}$ 的初值为 $\alpha_0=0$, 递推关系为
\[
    \alpha_{n+1}=\frac{4\alpha_n-3}{3\alpha_n+4}, \qquad (n\in\mathbb{N}).
\]
求 $(\alpha_n)_{n\geq 0}$ 的通项公式.
\begin{solution}

\end{solution}
\end{problem}
